% Latex-template for one-page abstract submissions.
% NUMTA2026 - Numerical Computations: Theory and Algorithms
% 21-27 June 2026, Falkensteiner Garden Calabria Resort, Curinga (Italy)

\documentclass[oribibl]{llncs}

% Please, do not use your own macros/redefinitions.

\usepackage[russian]{babel}

\usepackage{fancyhdr}
\renewcommand{\headrulewidth}{0pt}
\renewcommand{\footrulewidth}{0pt}

\begin{document}
\cleardoublepage

\title{Finding a Zero Value of a Non-negative Univariate Function}
\author{Konstantin Barkalov, Alexander Zaitsev, Roman Strongin}

%%% Please note that the first author is supposed to be the speaker.

\institute{Lobachevsky State University of Nizhni Novgorod, Address1, Country1 \\
\email{konstantin.barkalov@itmm.unn.ru, Name2@university2.edu}}

\maketitle

\thispagestyle{fancy}

\textbf{Keywords.} Keyword1; keyword2; keyword3.

 \vspace*{0.5cm}

%-----------------------------------------------------------------

В данной работе рассматриваются одномерные задачи липшицевой глобальной оптимизации и методы их решения. Важность рассматриваемых задач обусловлена следующими причинами.
Существует значительное количество практических приложений, в которых необходимо решать одномерные задачи глобальной оптимизации (см., например, [1]). Также важно разрабатывать и исследовать одномерные методы, поскольку их можно успешно обобщить на многомерный случай с помощью различных схем.

В целом, предположение о липшицевости целевой функции является типичным для многих подходов к разработке алгоритмов глобальной оптимизации (см, например, [2]). При этом надлежащая работа таких алгоритмов (т.е. сходимость к глобальному минимуму) обеспечивается либо при использовании достаточно хорошей оценки константы Липшица (задаваемой заранее или вычисляемой в процессе работы алгоритма), либо при использовании всего множества возможных значений константы.

В данной работе предлагается новый алгоритм, предназначенный для решения задач глобальной оптимизации, в которых значение целевой функции в точке минимума предполагается известным, но сама эта точка подлежит определению. Не ограничиваю общности, известное значение функции в точке глобального минимума можно полагать равным нулю. Задачу такого вида можно интерпретировать как задачу поиска параметров модели, обеспечивающих минимум невязки между расчетными и экспериментальными данными. Для ее решения, конечно же, можно применять известные методы липшицевой оптимизации. Однако в отличие от них, предлагаемый в данной работе алгоритм 1) гарантирует сходимость к глобальному минимуму целевой функции 2) не требует для своей работы ни априорной, ни апостериорной оценки значения константы Липшица.

\textbf{Acknowledgements.}

This research was supported by the following grants...

 \vspace{0.5cm}

\begin{thebibliography}{4}

\bibitem{Book} Cantor G. (1955) \emph{Contributions to the founding of the theory of transfinite numbers}. Dover Publications, New York.

\bibitem{Article} Gordon P. (2004) Numerical cognition without words: {E}vidence from {A}mazonia. \emph{Science}, Vol.~306,
pp.~496--499.

\bibitem{Proceedings} Author I., Author I.I. (2019) An abstract. In Proceedings of the \emph{International Conference``XXX''} (ed. by Editor E.), Place (Country), p.~1.

\bibitem{BookChapter} Author X. (2019) A chapter. In \emph{YYY Book} (ed. by Editor E.), Place (Country), pp.~10--20.

\end{thebibliography}

\end{document}
